\section{Methods}
\label{sec:methods}

\subsection{Setup and notation}
\label{sec:notation}

Our primary subject is \texttt{Qwen3-0.6B}, run locally; the cross-model study
(\Cref{app:cross-model}) extends to four further models. The data are the
\textsc{SESGO} items~\citep{robles2025sesgo}; each item is rendered into a prompt
$x$ that poses a question and lists answer options realizing the \emph{roles}
\[
  y \in \{\, \mathrm{target}\;(t),\ \mathrm{other}\;(o),\ \mathrm{unknown}\;(u) \,\}.
\]
The gold role $y^\star$ is set \emph{per condition}: on \emph{ambiguous} items
$y^\star = u$ (abstain); on \emph{disambiguated} items $y^\star$ is the labelled
role ($t$ or $o$). A \emph{scaffold} $s$ is a one-sentence debiasing preamble; we
write $x \oplus s$ for the prompt with $s$ prepended and compare $\LLM(x)$ against
$\LLM(x \oplus s)$. Throughout, $\mathbf{1}[\cdot]$ is the indicator and
$\delta_u$ is the one-hot distribution placing all mass on $u$.

We read every item at \emph{three} levels: two teacher-forced readouts (over all
three roles, and over just the two named roles $t,o$ with the unknown option
removed) and a greedy reasoning decode. These three readouts are shared by every
study; pipeline-specific variants and measures are defined in each appendix.

\paragraph{Teacher-forced readouts (three-option and two-option).}
We forbid reasoning (prefill a skip-thinking prefix and teacher-force each
option) and, for each role $y$, read the option-label token's conditional
log-probability at the answer position, $\ell_y = \log P(\text{label}_y \given
x)$, scored in one batched forward pass. The \emph{three-option} readout
renormalizes over all three roles,
\begin{equation}
  p^{\mathrm{nt}}_y \;=\; \frac{\exp(\ell_y)}{\sum_{y'} \exp(\ell_{y'})},
  \qquad y,y' \in \{t,o,u\},
  \label{eq:pnt}
\end{equation}
with chance $1/3$; the \emph{two-option} forced-choice readout drops the unknown
option and renormalizes over $\{t,o\}$ only (chance $1/2$), so it cannot abstain
and is undefined ($N/A$) on ambiguous items. Options are scored at \emph{display}
positions and remapped to roles through the item's position$\to$role permutation,
averaging out position bias. The prediction is $\hat{y} = \argmax_y p_y$ (ties
resolve to $u$ in the three-option case).

\paragraph{Greedy-thinking readout (decoded).}
We take a greedy reasoning decode (temperature $0$), strip the reasoning block,
and parse the committed role $r \in \{t,o,u,\varnothing\}$, where $\varnothing$
marks an undetectable or truncated draw (dropped); chance $1/3$.

\paragraph{Per-condition accuracy.}
The headline behavioral score is \emph{per-condition accuracy}---the fraction of
items whose prediction matches the condition's gold role $y^\star_q$:
\begin{equation}
  \mathrm{acc} \;=\; \frac{1}{|\mathcal{Q}|}\sum_{q\in\mathcal{Q}}
    \mathbf{1}[\, \hat{y}_q = y^\star_q \,].
  \label{eq:acc}
\end{equation}
On ambiguous items ($y^\star_q = u$) this is \emph{abstention accuracy}; on
disambiguated items it is the rate of naming the correct labelled role, and the
$t$-vs-$o$ accuracy gap within disambiguated items is the bias signal. We report
Wilson 95\% confidence intervals on every rate. The three-option role masses
$(p^{\mathrm{nt}}_t,p^{\mathrm{nt}}_o,p^{\mathrm{nt}}_u)$ give a companion view of
where probability leaks.

\subsection{The debiasing scaffold}

The scaffold $s$ is a one-sentence preamble prepended to the question to push the
model toward the $u$ answer; we compare four bilingual scaffolds against the
no-scaffold baseline (\Cref{tab:scaffolds}). Each scaffold uses a distinct
mechanism and is authored in matched English and Spanish (EN/ES). All four derive
from the structure-aware diversity methodology of
\citet{riossialer2026structureaware}.

\begin{table}[t]
  \centering
  \caption{The four bilingual debiasing scaffolds plus the no-scaffold baseline.
    Each scaffold pushes the model toward the \emph{unknown} answer by a distinct
    mechanism and is authored in matched English and Spanish (EN/ES) text.}
  \label{tab:scaffolds}
  \small
  \begin{tabular}{@{}p{0.40\linewidth} p{0.54\linewidth}@{}}
    \toprule
    \textbf{Scaffold id} & \textbf{Mechanism (one line; EN/ES matched text)} \\
    \midrule
    \texttt{(none)} & Baseline: the question with no preamble. \\
    \texttt{interpretive\_direction} & Reframes the task as judging whether the
      text identifies anyone---do not guess or fill gaps. \\
    \texttt{prior\_dominance\_warning} & Warns that models default to dominant /
      stereotyped assumptions under uncertainty and present them as fact. \\
    \texttt{insufficient\_information\_unknown} & Licenses ``unknown'' as the
      accurate answer when the text does not specify, not a failure to answer. \\
    \texttt{intent\_and\_register\_respect} & Demands the people be read on their
      own terms; attribute nothing the text does not state. \\
    \bottomrule
  \end{tabular}
\end{table}

\subsection{The five studies}

The framework comprises five studies (\Cref{tab:studies}). They share the item
pool and the readouts above, and differ in which surface variations are exercised,
which scaffolds are applied, and which readout is run.

\begin{table*}[t]
  \centering
  \caption{The five studies. Each row is one experiment, specifying the surface
    \emph{format variation} exercised, the \emph{models / scaffolds} applied, the
    \emph{readout(s)} run, and \emph{what is measured}. ``3-opt''/``2-opt'' are
    the three- and two-option teacher-forced readouts; ``greedy-thinking'' is the
    decoded reasoning readout.}
  \label{tab:studies}
  \small
  \setlength{\tabcolsep}{4pt}
  \renewcommand{\arraystretch}{1.15}
  \newcommand{\rr}[1]{\raggedright\arraybackslash #1}
  \begin{tabular}{@{}>{\rr}p{0.14\linewidth} >{\rr}p{0.18\linewidth} >{\rr}p{0.15\linewidth} >{\rr}p{0.23\linewidth} >{\rr}p{0.21\linewidth}@{}}
    \toprule
    \textbf{Study} & \textbf{Format variation} & \textbf{Models / scaffolds} &
      \textbf{Readout(s)} & \textbf{What is measured} \\
    \midrule
    \texttt{baseline} & none & 1 model, none & 3-opt $+$ 2-opt $+$ greedy-thinking &
      per-condition accuracy $+$ role mix, by bias category \\
    \texttt{cross\_model} & none & 5 models, none & 3-opt $+$ 2-opt $+$ greedy-thinking &
      accuracy vs.\ model size (abstention $+$ disambiguated) \\
    \texttt{stability} & $18$ ($3$ styles $\times$ $6$ perms) & 1 model, none &
      3-opt $+$ 2-opt & answer consistency $C_q$ across cosmetic variation \\
    \texttt{selection} & none & 1 model; all 4 scaffolds (Tab.~1) scored & 3-opt &
      per-scaffold abstention $\to$ select $s^\star$ \\
    \texttt{divergence} & none & 1 model, none & 3-opt $+$ 2-opt $+$ greedy- $+$ sampled-thinking &
      per-readout accuracy $+$ reasoning JS-div from unknown \\
    \texttt{geometry} & scaffold + none & 1 model, scaffold vs.\ none & greedy-thinking $+$ activations &
      silhouette separability of answer geometry by every label \\
    \bottomrule
  \end{tabular}
\end{table*}

In one line: \textbf{baseline} reads the full grid by bias category;
\textbf{cross\_model} repeats it across five models against model size;
\textbf{stability} measures answer consistency across cosmetic variation;
\textbf{selection} picks the debiasing scaffold $s^\star$ by per-scaffold
abstention accuracy (full-data, all four scaffolds; \Cref{app:fulldata});
\textbf{divergence} compares all four readouts and the reasoning distribution; and
\textbf{geometry} reads the answer-position representation. Each pipeline's
specific methodology, bulleted results, and figures are given in its own appendix
(\crefrange{app:baseline}{app:geometry}).

%%%%%%%%%%%%%%%%%%%%%%%%%%%%%%%%%%%%%%%%%%%%%%%%%%%%%%%%%%%%%%%%%%%%%%%%%%
